\subsubsection{Augmentasi Data}

Augmentasi diterapkan pada data pelatihan untuk menambah variasi data yang diterima model selama proses pembelajaran. Penerapan augmentasi bertujuan membantu model mempelajari variasi tampilan objek dan mengurangi kecenderungan model menghafal karakteristik citra tertentu \cite{yang2022}.

Augmentasi dilakukan secara dinamis melalui \textit{pipeline} pelatihan YOLO11-seg. Dengan mekanisme ini, transformasi dapat diberikan ketika data pelatihan dimuat tanpa harus menghasilkan salinan citra baru secara permanen. Transformasi pada citra juga diterapkan pada anotasi sehingga posisi \textit{segmentation mask} tetap sesuai dengan objek setelah augmentasi.

Jenis transformasi dipilih dengan mempertimbangkan karakteristik citra dan kondisi penggunaan sistem. Transformasi dapat mencakup perubahan posisi, skala, orientasi, serta karakteristik warna dan pencahayaan selama transformasi tersebut tidak mengubah identitas utama objek atau menghasilkan kondisi yang tidak relevan terhadap penggunaan sistem.

Augmentasi acak tidak diterapkan pada data validasi dan pengujian agar kedua bagian tersebut tetap digunakan sebagai data evaluasi yang merepresentasikan citra aslinya. Konfigurasi augmentasi yang benar-benar digunakan pada proses pelatihan dicatat pada tahap pembuatan sistem.
